\section{白書は対象年と集計範囲を先に読む}
\label{theory:env:climate}
\label{theory:env:oil2012}
\label{theory:env:energy2010}
\label{theory:env:energycapita}
\label{theory:env:energymix2015}
白書の発行年、統計の対象年、政策の策定年、目標年は違う。過去問の年を現在の統計へ置き換えると、問の意味まで変わる。数量の輸入依存度、金額の輸入比率、人口当たり一次エネルギー、電源別発電量は分母を別々に固定する。部門別CO\(_2\)でも電力を需要部門へ配分するかで比較が変わる。

温室効果ガスのCO\(_2\)換算はガスの質量をそのまま足す操作ではない。気候の条約目的と各議定書の数量約束も区別する。2015年に策定した2030年度見通しと、後年に改訂した同じ2030年度の比率は別資料である。目標値は実測結果ではない。歴史表の追加照合が必要な問題はその旨を個別に表示する。\cite{env58_ghg2008,env58_oil,env58_mix2015}

\section{活動量に排出係数を掛けてから人数で割る}
\label{theory:env:co2}
\label{theory:env:lng}
\label{theory:env:heating}
\label{theory:env:resource}
電力 \(P\) に時間 \(t\) を掛けて電力量を得る。排出係数 \(e\) を使うと \(M=Pt e\)。WとkW、時間と日数をそろえ、一人当たりなら最後に利用人数で割る。車は走行距離を燃費km/Lで割って燃料使用量を求める。係数は出題条件を使い、現行の電力係数へ置換しない。

同じ質量の気液体積比は \(V_l/V_g=\rho_g/\rho_l\)。理想気体で標準状態を指定されたメタンなら、分子量とモル体積から気体密度を独立に作れる。燃料発熱量はMJ/kgとMJ/L、標準立方メートル当たりを混ぜない。資源の存在量と経済的可採量も違う。太陽光の概算では受光面積、日照、変換効率、年間時間を含め、定格電力を一年中の出力とはしない。\cite{env58_heating}

\section{循環利用の優先順位と役割分担}
\label{theory:env:wasteorder}
\label{theory:env:recyclelaw}
\label{theory:env:wasteterm}
\label{theory:env:balance}
廃棄物の発生抑制を先に考え、再使用、再生利用、熱回収、残るものの適正処分へ進む。環境負荷等による合理的な留保もある。再使用は形を保つ利用、再生利用は材料へ戻す利用なので同じ語にしない。排出者・自治体・小売・製造者が担う回収・収集・再商品化の役割と費用は、制度ごとに確認する。\cite{env_circular_law,env_appliance,env58_wastegoal}

フロー図では各節点の入力と出力が保存される。排出総量は直接再利用、中間処理、直接処分の和、中間処理は処理後再利用、処理後処分、減量の和となる。減量を再利用へ二重に数えない。RDFは加工燃料、E-wasteは電気電子機器の廃物を指し、環境に良いという判定を定義へ付け加えない。

\section{環境管理手法は目的と単位で区別する}
\label{theory:env:lca}
\label{theory:env:accounting}
LCAは製品等の機能単位とシステム境界を定めて全過程の環境影響を評価する。CFPは温室効果ガスのCO\(_2\)換算量に着目する。報告は利害関係者への情報伝達、監査は環境管理の実施と適合性を点検する。環境会計は環境保全コストと効果を貨幣又は物量単位で認識・測定・伝達する。\cite{env58_accounting}

排出や吸収を通じて気候変動の原因を減らす緩和と、影響に備えて被害を減らす適応を区別する。環境基準は環境の望ましい状態を示す。個別排出源の規制値とは対象と責任が違い、基準超過の測定だけで任意の一企業の法違反を確定しない。\cite{env58_pmstandard}

\section{汚染対策は生成・移動・除去の工程に分ける}
\label{theory:env:treatment}
\label{theory:env:air}
PM2.5は概ね粒径2.5\(\,\mu\)m以下の粒子の集合で、2.5付近一種類ではない。告示の測定定義は分粒装置の50\%分離性能を含む。SO\(_2\)は硫黄を含む燃料、COは不完全燃焼、光化学オキシダントは前駆物質と日射による二次生成という経路の違いがある。\cite{env58_pmstandard,env58_pm25}

下水の一次処理は沈殿等、二次処理は活性汚泥等の生物作用、高度処理は除去対象に応じた追加工程である。処分場の遮水と浸出水処理、緑化の熱移動、焼却の完全燃焼と排ガス再合成防止も、どの工程を制御するかで整理する。効果は濃度、温度、滞留時間等に依存し、名称だけで全汚染を除去する保証とはしない。

\section{再エネ導入と発電・貯蔵の仕組み}
\label{theory:env:fit}
\label{theory:env:generation}
2012年開始のFITは対象電源、調達価格、接続、賦課金、費用調整を組み合わせた制度である。原則義務の記述に正当な拒否理由がある場合、例外を消して「一切拒否できない」へ変えない。出題後の制度改正は別に扱う。\cite{env58_fit2013}

コンバインドサイクルは高温の燃焼ガスでガスタービン、その排熱で蒸気を作り蒸気タービンへ進む。燃料電池は電気化学、電気二重層キャパシタは界面の電荷蓄積、二次電池は可逆な化学反応を利用する。揚水は位置エネルギーへ変換して貯蔵するので、発電機効率と充放電の往復効率を分ける。

\section{環境法の目的と生物多様性}
\label{theory:env:warminglaw}
\label{theory:env:biodiversity}
温暖化対策では排出抑制に加え吸収作用の保全・強化も見る。対象ガスの追加、報告義務の規模、計画の検討周期は時点の法令を確認する。現行の目標年や条番号を過去の条文へ混入させない。

生物多様性の保全、構成要素の持続可能な利用、遺伝資源の利益の公正衡平配分は異なる三目的である。組換え生物の扱いと利益配分の議定書も別である。外来種の防除と動物の扱いへの配慮は同時に検討でき、一律駆除禁止又は配慮不要という二択にしない。

\section{知的財産と専門職の責任}
\label{theory:env:ip}
\label{theory:env:ethics}
特許発明は自然法則を利用した高度な技術的思想。知的財産全体には著作物・商標・営業秘密等も含む。保護対象、権利取得、国の施策責務を別々に読む。\cite{ip_patent_statute,ip_broad_scope}

専門職は力量の範囲と公益に関する高い責任を持つ。職務規定に項目がないだけで責任がなくなるとはいえない。能力外の仕事には、必要な力量を持つ者との協働・委任・引継ぎなどを考え、不十分な能力のまま安全判断を引き受けない。\cite{ipej_ethics2011}

\section{技術史は発明・改良・発表・商業化を分ける}
\label{theory:env:historypairs}
\label{theory:env:historyorder}
\label{theory:env:industrial}
人物名と代表的業績の対応を確認した後、世紀と発表・発明の順を整理する。Becquerelの放射能とキュリー夫妻の新元素、NewcomenとWattの動力機械、ナイロンの合成と商業化はそれぞれ異なる段階である。ノーベル賞の授賞年をそのまま発見年へ移さない。\cite{env58_curie,env58_steam,env58_nylon}

産業革命は発明家・職人・事業者の改良、生産機械、工場制度と労働の変化を含む。大学との接点がある個別発明から当時の全技術が大学主導だったと一般化しない。原典の「発見」や「フィルム乾板」のような粗い表現は、厳密な史料年表と区別して留保を残す。

\section{科学と社会のコミュニケーション}
\label{theory:env:communication}
専門情報を発信するだけでなく、非専門家の疑問、期待、価値判断を理解し研究者側へ伝える双方向の対話を考える。基礎科学と応用技術の内部連携や、媒体制作に科学技術を使うことは、この目的の全体を代替しない。伝え手側も社会への理解を深める。\cite{env58_communication}
